\section{Constant Step Aggregation of Plonkish with zk-cq}
  For the ZkregRegPlus Project, we need a zksnark proof system
which supports the following features:
\begin{enumerate}
	\item A commit-and-prove zkSNARK system which allows one
	to commit to some witnesses, and then some other witness
	are generated using Fiat-Shamir.
	\item Support lookup argument with the number of lookups (quries)
	much less than the lookup table.
	\item Suppport multi-instance folding. We use this mainly as
	an alternative to proof-aggregation (such as \cite{Bunz21,SNARKPack})
	as folding provides much faster prover concrete performance.
	The number of folding steps, as shown in Section 2, is no more
	than $100$. Hence, we do not need a 2-cycle IVC scheme where
	the IVC verifier is embedded in the proof itself.
\end{enumerate}

Given the above, we ends up with a design of a $3$-degree
Plonk system \cite{GWC19} (which natively allows Fiat-Shamir of
multiple-stage inputs), plus the adoption of QA-NIZK for commit
and prove so that the zk-cq protocol can be applied for lookup.
We then apply folding as an alternative for proof aggregation,
by leveraging commit-and-prove (also for folding lookup argument).
This results in a system that needs only one pairing friendly curve,
with prover complexity $O(kn\log(n))$ and verifier
complexity $O(\log(n))$ where $k$ are the number of instances
and $n$ is the circuit size for one circuit. But the memory
footprint is $O(n)$ for one circuit (thanks to folding).

We make a slight adaptation of the Plonkup relation in
ProtoStar. In our structure, a lookup table can be a concatenation
of many smaller tables, so a lookup table is a 2-column table
with the first column indicates the sub-table ID.
\begin{definition} \cite{ProtoStar}[Definition 15] A {\em Plonkup Relation}
is a tuple $\caC = (n, T, \sigma; c, d, [s_i, G_i]_{i=1}^m; L, \myvec{t})$
s.t. $\delta: [cn] \rightarrow [cn]$ is a permulation,
$(c,d,[s_i, G_i]_{i=1}^m)$ are custom gates, and $L [cn] \rightarrow \F$
maps a variable index to the corresponding sub-table ID (if it's $0$
it means that no lookup needed for this variable). The vector $w$
can be written as $\myvec{w} = io \concat \myvec{u} \concat p \concat \myvec{v}$
where $p$ is the Fiat-Shamir random, defined as
$\hash(io \concat \myvec{u})$, and $\myvec{v}$ are rest of the
witness values (which might depend on $p$). Note that
halo2 implementation of Plonup system already provides commit-and-prove
Fiat-Shamir.
\end{definition}

By extending the ProtoGalaxy framework \cite{ProtoGalaxy},
Plonkup relation $\caC$ can always be compiled into a
relation below. The difference is the addition of $\myvec{t}$
(as we do not plan to encode $\myvec{t}$ into the Plonk relation).

\begin{definition} A ProtoGalaxy interactive protocol for
$\caC$ is defined as a tuple $(k, d, n, \myvec{w}, L, \myvec{t})$
$\myvec{t}$ is a pre-defined lookup table
where transcript $\myvec{w}$ is a concanetation of
$k$ rounds of message vectors $\myvec{m}_i$ and 
Fiat-Shamir random nonce $r_i$, and a $d$-degree mapping
$f: \F^M \rightarrow \F^n$ s.t. 
$f(\myvec{w}) = 0^n$. In addition, $L$ is a collection of
indexes that map $\myvec{w}$ to a subset $\myvec{v}$ s.t.
$\myvec{v} \lookup \myvec{t}$. Note that $L$ is fixed by circuit,
but witness vector $\myvec{w}$ is per instance.
\end{definition}

\begin{definition} Given $\beta$ define 
$\myvec{\beta} = (\beta^{2^0}, \beta^{2^1}, ..., \beta^{2^t}$
where $t= \log(n)$, and $pow_i(X_0, ..., X_{t-1}) = \prod_{\ell \in S} X_\ell$
for $i-1 = \sum_{j\in S} 2^j$. Define
a committed instance of $\caC$ as
$(\phi, \Commit_t ;\omega)$ where $\phi = \cm(\omega)$
and $\Commit_t = \cm(L(\omega))$.  Here the $\Commit_t$ is the
extra commitment to the subset of $\omega$ which needs to be
checked with lookup.
\end{definition}

\begin{definition} A randomlized relaxed version is defined as
$((\phi,$ $\Commit_t, \myvec{beta}, e); \omega))$ s.t.
$\phi = \cm(\omega)$, $\sum_{i\in [n]} pow_i(\myvec{\beta})f_i(\omega) = e$,
and the subtable in $\Commit_t$ contained in lookup table.
\end{definition}

We now adapt the ProtoGalaxy scheme with GIPA so that
We can (a) connect ZVC inputs and outputs, (b) batch proof,
(c) extract commitments of query table for look-up.

Recall that the ProtoGalaxy verifier outputs the following:
\[
	\phi^* = L_0(\gamma) \phi + \sum_{i\in [k]}L_i(\gamma) \phi_i
\]
And, the prover outputs the following:
\[
	\omega^* = L_0(\gamma) \omega + \sum_{i\in [k]}L_i(\gamma) \omega_i
\]
We need to address the committed query tables $\myvec{t}_i$ in
each $\omega_i$ (and the $\Commit_t$ for $\omega$). 
We compress them as follows:
\begin{equation}
	\Commit_t^* = L_0(\gamma) \Commit_t + \sum_{i\in[k]} L_i(\gamma) \Commit_{t_i} \label{eq-tstar}
\end{equation}

Note that we do not want the verifier to view the entire
$\phi_i$, and also $\Commit_{t_i}$. So the prover can just
provide an AFGHO commitment to $\{\Commit_{t_i}\}_{i=1}^k$,
and run MIPP argument to prove the validity of
\ref{eq-tstar}.

{\bf Problem:} if all info are packed, there is no way for
verifier to check each proof $\phi_i$ is H-consistant,
that is in each $\phi_i$: $r_i = \hash(\Commit_i, r_{i-1})$.
This seems hard to address? But {\em if there is only one round,
$k$ = 1, then there is no need to compute this round hash? - 
this in Plonkup has only one round for plonkup, excluding
the handling of lookup relation!} So the key idea here is to limit
the plonish system to 1 round (without verifier challenge).

A similar problem arises that for multi-set equivalence test,
there is a verifier challenge needed. This is avoided by
letting the prover committing to all stage-1 messages as one
commitment and generate one ``global" verifier challenge based
on the global commitment. 

Also how do verifier check that $\Commit_{t_i} = \Commit_{L(\phi_i)}$
-- this can be resolved because the prover can run a QA-NIZK proof
and provides the proofs and these proofs can be verified
using zk-TIPA proof.

In the following we briefly describe the scheme below (later
{\bf TODO:  formalize it}.

\begin{enumerate}
  \item The prover has $k$ instances of Plonkish relation. Do not
  generate the intermediate wires yet, let $\myvec{u}^{(i)}$ be the
  stage-1 input witness to commit for instance $i$. 
  Let $\myvec{u} = (\myvec{u}^{(i)})_{i=1}^k$. Compute
  $\GCommit_{\myvec{u}}$. (decision to make: AFGHO commitment
  or Perdersen vector to the entire vector).

  \item The prover then run Fiat-Shamir and generates
  the verifier challenge $r$ over $\GCommit_{\myvec{u}}$.
  The prover then populates the rest of wire values for
  the $k$ instances of Plonkish relation.
  
  \item The prover then translates all $k$ Plonkish instances
  to ProtoGalaxy form: $(\phi_i, \Commit_{t_i}); \omega_i)$.
  Note that to avoid verifier complexity to be linear over $k$,
  the prover computes $\GCommit_\phi = \gcm(\{\phi_i\}_{i=1}^k)$.
  and similarly $\GCommit_{\Commit_{t}} = \gcm(\{\Commit_{t_i}\}_{i=1}^k)$.
  In the ProtoGalaxy protocol, when needs to send over
  $F_1, ..., F_t$ and $K_1, ..., K_t$ (co-efficients of polynomials
  $F(X)$ and $K(X)$, the prover sends their commitment
  $\Commit_F$ and $\Commit_K$ instead. Then the verifier computes
  Fiat-Shamir challenge $\sigma$ and $\gamma$ over these commitments
  instead.

  	An improvement we can make over ProtoGalaxy is that
	in the protocol, all prover messages can be similarly packed
	as a commitment and the corresponding proof. We also want to
	avoid linear size verifier work, e.g., computing
	$\myvec{\sigma} = \{\sigma^1, \sigma^2, ..., \sigma^{2^t}\}$.
	But instead, the prover will provide a commitment to it and
	the corresponding proof for the wellformned ness of such commitment.

	The last step of verifier computation of the following can
	be improved similarly:
	\[
		\phi^* = L_0(\gamma) + \sum_{i\in [k]} L_i(\gamma) \phi_i
	\]
	The prover provides a commitment to the vector of $L_i(\gamma)$
	given $\gamma$ and provides a validity proof. Notice that
	$L_j(x) = L(x) \frac{\mu_j}{x-\mu_j}$ where $\mu_j$ is the
	$\mu^j$ for root of unity $\mu$, and $L(x) = \prod_{i\in[k]} (x-\mu_i)
	 = x^k -1$. Such proof can be very conveniently captured by
	 a plonkish circuit, and the proof generated and commitments
	 extracted via column commitments. Given given
	 the commitment to $\{L_i(\gamma)\}_{i=1}^k$, the prover 
	 can run an MIPP proof to prove the validity of the above formula
	 given all $\phi_i$'s are contained in AFGHO commitment 
	 $\GCommit_{\phi}$.

  \item Now the prover has the aggregated instance
  $( (\phi^*, \Commit_t^*); \omega)$. It translates back the instance
  to a Plonkish system and feed it to the Plonk and generate the proof.
  Then it needs to establish the equivalence between the aggregated
  ProtoGalaxy instance and the generated Plonkish instance.
  This is constructed by a QA-NIZK to prove the equivalence of
  witness columns. We also need to justify the Fiat-Shamir is
  really in the $\omega$, we can make it a public input or use
  vector commitment (make it public a wise choice?)

  \item Finally, similarly run a TIPP proof for vector of $\Commit_t$ 
  and a valid global $\Commit$ for all queries. Then run a separate
  zk-CQ proof over it.

\end{enumerate}

{\bf Summary:} Prover cost will be $O(kn)$ for folding,
$O(n\log(n))$ for the last step of proving and
$O(k\log(k))$ for TIPP/MIPP proofs. The verifier cost
will be $O(\log(k))$. Compared with the 
aggregation scheme directly with GIPA, the concrete
cost of prover is faster $2kn$ MSMs vs e.g., $22kn$ MSMs.
The verifier also has much smaller concrete cost as
there is only one 1 MIPP and 2 TIPP proofs (instead of
running TIPP proof for every single pairing operation
in the proof).

{\bf Translation Between Systems:} We have to translate
between (1) Plonkish system (2) ProtoGalaxy encoding.

From Plonkish to ProtoGalaxy: The main job is to 
define the degree $d$ function $f(\myvec{w})$ that
performs the verifier work. It basically consists of
two checks: (1) the permutation check which is degree 1
(2) those custom gates which can be up to degree 3.
Basically, the verifier function $f(\myvec{w})$
can be expressed as $O(|\myvec{w}|)$ degree 3 constraints,
more exactly $O(n)$ constraints $f_i(\myvec{w})$.
Since $f_i$ is up to degree 3, the input parameter
can be represented compactly. Now, the $f_i(\myvec{w})$
will be involved in computation in ProtoGalaxy:
\begin{enumerate}
	\item Computing $F(X)$, since $\myvec{w}$ is already given,
	$f_i(\myvec{w})$ can be computed in constant $O(1)$ time,
	then $F(X)$ can be computed in $O(n)$ time.

	\item Computing $G(X)$: use the FFT approach at least 
	could get us $O(n\log(n))$ complexity. The total needs
	$dknC$ as in Claim 4.5.
\end{enumerate}

From ProtoGalaxy to Plonkish: The prover has $\myvec{w}$
and the set of $f_i$ specification, and tries to prove
two things: (1) $\phi = \commit(\myvec{w})$. (2)
$\sum_{i\in [n]} \pow_i(\beta) f_i(\omega) = e$. The idea
of handling is described as below:
\begin{enumerate}
	\item $\phi = \commit(\myvec{w})$ can be simply
	used to extract from the Plonkish proof commitment.
	\item The logic of $\sum_{i\in [n]} \pow_i(\beta) f_i(\omega)$
	can be simply encoded as a plonkish system with
	each $f_i$ encoded one by one using standard
	arithmetic operations.
\end{enumerate}



